\subsection{Exploratory screening for highly similar but non-identical recordings}
\label{sec:p2-whole-record-near-overlap}
The exact-identity audit establishes 49 groups of
\emph{byte-equivalent processed recording contents}
for the primary analysis; it cannot identify
near-duplicates produced by scaling, quantisation,
partial corruption or changes in physical calibration.
We therefore performed a separate retrospective
screen of all $\binom{50}{2}=1,225$ released-ID pairs.
Using the same pinned 8--30~Hz/160~Hz processed
40-trial arrays, each complete trial-aligned
record was vectorised, mean-centred and normalised
to unit norm. Its full-record Pearson/cosine
similarity was compared to precommitted
diagnostic cutoffs 0.95 and 0.9999.
Event-label sequence agreement was checked
separately. This screen was developed after
the known ID4/5 finding and is \emph{exploratory}.
It is not invariant to trial reordering, time
shifts, altered events or channel permutations.

Precisely three of the 1,225 pairs exceeded
both the 0.95 and 0.9999 thresholds.
The complete processed contents of ID4/5
were exactly identical as established earlier,
and their normalised similarity was 1.0.
Two \emph{additional} pairs exhibited
near-unity whole-record similarity:
ID47/49 had $0.99999999166$, and ID11/15
had $0.99999998340$. All three pairs
had identical aligned 40-trial label sequences.
The remaining 1,222 pairs fell below 0.95
under this particular aligned-record measure.
These new pairwise comparisons were
executed on the same public release using
the frozen screen and saved at
\url{https://github.com/lindicaphxag-tech/kaggle/actions/runs/38321951859}.

Importantly, the original released ID47/49
pair was previously used as a
\emph{negative control for exact identity}:
a separate original EDF digital-payload audit
confirmed byte-for-byte equality in only one of
33 raw channels. It is \emph{not} a sound
negative control for waveform near-overlap.
A lack of raw digital byte identity can coexist
with nearly identical processed numerical
waveforms. We therefore retain the primary
49 \emph{exact-content} groups but
\emph{do not} infer there are 49 fully
independent electrophysiological recordings
under any tolerance- or acquisition-based
notion of identity. Likewise, we do not
prematurely merge 47/49 or 11/15 into
single people or claim an effective 47-patient
sample. Such interpretation requires inspecting
the original acquisition/release mechanism
and comparing calibrated raw/processed signals
under matched temporal/channel alignment.
Whole-record correlation is only a trigger
for this technical forensic inquiry.

The related high-similarity result does not
alter the original 49 exact-group
$2\times2$ numerical factor analysis,
whose estimand is conditional on \emph{exact}
processed-content identity. It narrows the
interpretation of that analysis: exact
deduplication alone is not sufficient
evidence of physiologically independent
recordings in this public release.
